% !PRINTED none
\gesperrt{Twentieth Chapter.}

The Multiplier Equation in Complex
Multiplication.

\begin{center}\rule{3cm}{0.4pt}\end{center}

§ 134. \gesperrt{The Class Invariant $\gamma_3(\omega)$.}

We have now to test the function
\[
\gamma_3(\omega) = \sqrt{j(\omega) - 12^3}
\]
for its property as a class invariant.

We have first, for a variable $\omega$, the multiplier equation
 of the 1st level (§ 72, 3), according to which, if
\weq{1}{n \equiv 3, \quad c \equiv 0\ (\mathrm{mod}\ 4),}
the $\nu$ quantity, % [sic] singular in the printed source
\weq{2}{M = (-1)^{\frac{a-1}{2}} \partial^3\left(\frac{\eta\left(\frac{c+\partial\omega}{a}\right)}{\eta(\omega)}\right)^6 \gamma_3(\omega),}
are merely permuted among one another by the linear
transformations. If we therefore set
\weq{3}{u = j(\omega), \quad v = j\left(\frac{c+\partial\omega}{a}\right),}
then between $u, v$ there subsists the invariant equation
\weq{4}{F(v,u) = 0,}
and we have
\[
F(x,u) = \Pi(x-v),
\]
when the product sign extends over the $\nu$ systems of values of
$a, b, c$. The sum
\weq{5}{F(x,u)\sum \frac{M}{x-v} = \psi(x,u)}
is, for an indeterminate $x$, unchanged by linear transformation,
 and therefore a rational function of $u$. But since the function
 has, for every finite value of $j(\omega)$, a finite
%CONT>>
